核电荷为 \(+Ze\)（\(Z\ge1\)，\(e>0\)）的类氢离子只含一个电子，考虑有限核质量造成的约化质量 \(\mu>0\)，真空介电常数 \(\epsilon_0\)。电子—核库仑势为 \(U(r)=-Ze^2/(4\pi\epsilon_0r)\)。已知电子基态的空间波函数形式为 \(\psi_{1s}(r)=C e^{-r/a_Z}\)，其中 \(a_Z=4\pi\epsilon_0\hbar^2/(\mu Ze^2)\)。
(1) 求归一化常数 \(C\) 与电子的径向概率密度，并计算最概然半径；(2) 从径向积分求 \(\langle r\rangle,\langle r^2\rangle,\langle1/r\rangle\)；(3) 求基态束缚能、库仑势能期望及动能期望，并验证 Coulomb 势的维里关系；(4) 求电子出现在半径 \(R>0\) 球内的累计概率，讨论核电荷数 \(Z\) 增大时轨道尺度和束缚能的标度。
